\documentclass[11pt]{article}

\usepackage[a4paper,margin=25mm,headheight=15pt]{geometry}
\usepackage[T1]{fontenc}
\usepackage[utf8]{inputenc}
\usepackage{lmodern}
\usepackage{microtype}
\usepackage{amsmath}
\usepackage{graphicx}
\usepackage{xcolor}
\usepackage{fancyhdr}
\usepackage[hidelinks]{hyperref}

\definecolor{Accent}{HTML}{315B7D}
\definecolor{Muted}{HTML}{5B6570}

\hypersetup{pdftitle={Supplementary Figures}, pdfauthor={Poikela et al.}}

\pagestyle{fancy}\fancyhf{}
\lhead{\small\color{Muted}Supplementary Figures}
\rhead{\small\color{Muted}Poikela et al.}
\cfoot{\small\thepage}
\renewcommand{\headrulewidth}{0.4pt}
\setlength{\parindent}{0pt}\setlength{\parskip}{0.6em}

\renewcommand{\thefigure}{S\arabic{figure}}
\newcommand{\Faq}{\textit{F.~aquilonia}}
\newcommand{\Fpo}{\textit{F.~polyctena}}
\newcommand{\ldw}{\ensuremath{\mathrm{ld}_w}}

\title{\vspace{-3em}\textbf{\color{Accent}Supplementary Figures}\\[0.3em]
  \Large LD decay, complexity reduction, ancestry sorting, architecture, and climate}
\author{Poikela et al.}
\date{}

\begin{document}
\maketitle
\vspace{-1.5em}
\begin{center}\small\color{Muted}
The analytical workflow is shown in the Materials and Methods. Figures are ordered
following the Methods sections. All analyses are descriptive and conditional; see the
Methods for interpretation.
\end{center}

% ---- LD decay and local LD support -------------------------------------
\begin{figure}[htbp]\centering
\includegraphics[width=0.86\textwidth]{figures/figS02_ld_decay.png}
\caption{\textbf{LD decay and local LD support.} Genome-wide linkage-disequilibrium
decay and the per-marker local LD support statistic \ldw{}, computed as the median
\(r^2\) between a marker and its neighbours within the \(\rho=0.95\) window. \ldw{} is
measured directly from the genotypes at marker resolution and is used to flag clusters
for reconsideration during complexity reduction.}
\label{figS:lddecay}
\end{figure}

\begin{figure}[htbp]\centering
\includegraphics[width=0.9\textwidth]{figures/figS03_roc_low_recomb.png}
\caption{\textbf{Discrimination of low-recombination regions.} Receiver-operating
characteristic curves for the local LD support \ldw{} and the window-level decay rate
\(a\), evaluated against map-based recombination rate at three definitions of low
recombination (lowest 5\%, 10\%, 25\% of windows). The area under the curve is
0.91--0.98 for \(a\) and 0.73--0.86 for \ldw{}; the decay rate is the stronger
window-level indicator, whereas \ldw{} supplies the marker-resolution measure.}
\label{figS:roc}
\end{figure}

\begin{figure}[htbp]\centering
\includegraphics[width=\textwidth]{figures/figS04_ld_tracks.pdf}
\caption{\textbf{Local LD support along example chromosomes.} Per-marker \ldw{} (grey)
along Chr26 and Chr10, each with a pronounced low-recombination region, shown against
the window decay rate \(a\) (left) and map-based recombination rate (right); all series
are min--max scaled. \ldw{} is elevated exactly where both \(a\) and recombination are
reduced.}
\label{figS:tracks}
\end{figure}

% ---- complexity reduction ----------------------------------------------
\begin{figure}[htbp]\centering
\includegraphics[width=0.82\textwidth]{figures/figS05_stage1_vs_combined.png}
\caption{\textbf{Two-stage complexity reduction.} Stage 1 clustering compared with the
combined two-stage result for Chr26, restricted to markers with high local LD support.
Fragmented Stage-1 clusters within each low-recombination region are consolidated into
single coherent clusters, the intended effect of the Stage 2 quality-gated merge.}
\label{figS:stage}
\end{figure}

\begin{figure}[htbp]\centering
\includegraphics[width=0.9\textwidth]{figures/figS06_fidelity_hist.png}
\caption{\textbf{Consensus fidelity.} Distribution of consensus round-trip fidelity
(\(\mathrm{score}_{\mathrm{eMLG}}=\mathrm{cor}\{\mathrm{round}(x),x\}^2\)) across five
\ldw{} flagging thresholds, separating single-marker clusters (no merge possible) from
genuine multi-marker merges. The separation clarifies that mass near the upper limit at
permissive thresholds reflects isolated markers rather than merges stopping early.}
\label{figS:fidelity}
\end{figure}

% ---- architecture -------------------------------------------------------
% (The extent and direction of sorting are shown in main-text Figure~1A,B; sort-class
% counts in Table~S1 and size-stratified aggregation in Table~S3.)
\begin{figure}[htbp]\centering
\includegraphics[width=0.78\textwidth]{figures/figS08_architecture.pdf}
\caption{\textbf{Genomic architecture of sorting.}
(A) Standardised relative differentiation (\(F_{ST}\)), between-parent allelic
difference (\(d_{xy}\)) and parental expected heterozygosity (\(\pi\)) across
recombination deciles: the differentiation signature is confined to the
lowest-recombination decile. (B) Direction of sorting versus recombination at the
LD-reduced-unit level, showing no strong recombination gradient. (The abundance
comparison of SNP- versus unit-level sorting is shown in main-text Figure~1C.)}
\label{figS:arch}
\end{figure}

% (Consensus-versus-representative validation, ~99% direction agreement, is reported
%  numerically in the Materials and Methods rather than as a figure.)

% ---- climate ------------------------------------------------------------
\begin{figure}[htbp]\centering
\includegraphics[width=0.8\textwidth]{figures/figS10_ancestry_confound.pdf}
\caption{\textbf{Ancestry--climate association.} Per-population mean \Faq{} ancestry
against the climate axes PC1 and PC2. PC1 is weakly correlated with ancestry
(\(\rho=-0.23\)); PC2 shows a moderate positive correlation (\(\rho=0.57\)) strongly
influenced by {\AA}land, motivating the use of the covariance matrix \(\Omega\) and the
{\AA}land-excluded configuration as primary.}
\label{figS:confound}
\end{figure}

\begin{figure}[htbp]\centering
\includegraphics[width=0.86\textwidth]{figures/figS11_dose_response.pdf}
\caption{\textbf{Diagnostic content versus climate association.} Proportion of
ancestry-diagnostic members against the cluster climate-significance rate, in
size-comparable strata (clusters with \(\geq20\) and \(\geq50\) markers), for PC1 and
PC2 in the primary configuration. Point size scales with the number of clusters; the
dashed line is the genome-wide baseline.}
\label{figS:dose}
\end{figure}

\begin{figure}[htbp]\centering
\includegraphics[width=0.86\textwidth]{figures/figS12_maf_power.pdf}
\caption{\textbf{Allele-frequency and differentiation sensitivity of the enrichment.}
The diagnostic-enrichment coefficient (percentage points per +10 percentage-point
climate rate) across nested models adding hybrid MAF, recombination, and BayPass XtX
(M0--M3), for PC1 and PC2. Blue intervals are model-based 95\% intervals; the dark M3
interval is the chromosome block bootstrap. The coefficient roughly halves and loses
significance once differentiation (XtX) is included: PC1 M0--M3 \(=+4.46,\,+3.48,\,
+3.44,\,+1.82\) (\(P=8\times10^{-4}\) to \(0.16\)); PC2 \(=+3.72,\,+2.46,\,+2.77,\,+1.69\)
(\(P=0.013\) to \(0.23\)). The M3 chromosome block-bootstrap intervals are
\([-3.9,\,10.0]\) (PC1) and \([-2.5,\,15.5]\) (PC2).}
\label{figS:mafxtx}
\end{figure}

\begin{figure}[htbp]\centering
\includegraphics[width=0.8\textwidth]{figures/figS13a_manhattan_PC1.png}\\[1mm]
\includegraphics[width=0.8\textwidth]{figures/figS13b_manhattan_PC2.png}
\caption{\textbf{Cluster-level climate association along the genome.} Bayes factors for
association with PC1 (top) and PC2 (bottom) in the primary configuration
({\AA}land-excluded, fixed-\(\Omega\)), with outlier clusters shown over the genome-wide
background.}
\label{figS:manhattan}
\end{figure}

% ---- dependence-aware uncertainty --------------------------------------
\begin{figure}[htbp]\centering
\includegraphics[width=0.72\textwidth]{figures/figS14_block_bootstrap.pdf}
\caption{\textbf{Dependence-aware uncertainty.} Genomic block-bootstrap 95\% intervals
(10,000 replicates) for the five central coefficients, from whole-chromosome and 10-cM
blocks, against the model-based intervals. The three sorting/architecture coefficients
exclude zero under both block definitions; the two climate coefficients do not.
Estimate [chromosome 95\%]: DI\,$\to$\,direction \(1.46\) \([1.41,1.52]\);
recombination\,$\to$\,direction \(-0.091\) \([-0.145,-0.037]\);
recombination\,$\to$\,magnitude \(0.052\) \([0.046,0.057]\);
climate\,$\to$\,diagnostic PC1 \(4.46\) \([-4.0,13.1]\) and PC2 \(3.72\)
\([-0.7,17.8]\). Every reconstructed model reproduced its full-data estimate before
resampling.}
\label{figS:boot}
\end{figure}

\end{document}
